\section{Literature search approach and critical synthesis}
\label{sec:lit}

\paragraph{Search approach.}
I started from the reproduced paper, LAION-CLAP~\cite{wu2023laionclap}, and searched in two directions.
\emph{Backward:} every system that LAION-CLAP compares against in its ESC-50 zero-shot table
(Wav2CLIP~\cite{wu2022wav2clip}, AudioCLIP~\cite{guzhov2022audioclip}, Microsoft CLAP~\cite{elizalde2023clap})
and the image-text model all of them build on, CLIP~\cite{radford2021clip}.
\emph{Forward:} Google Scholar and arXiv queries for papers that adapt a frozen language-audio model with a few labelled clips
(``CLAP few-shot'', ``training-free adapter audio-language'', ``Tip-Adapter audio''). This led to the
Treff adapter~\cite{liang2023treff}, which in turn uses Tip-Adapter~\cite{zhang2022tip} as its main baseline, so Tip-Adapter was added.
\emph{Recent work (2024--2026):} a second round of arXiv and Google Scholar queries restricted to 2024 onwards
(``audio-language model noise robustness'', ``CLAP test-time adaptation'', ``few-shot prompt learning audio language model'')
added four papers that build on CLAP-style models: a competing labelled few-shot method (PALM~\cite{hanif2024palm}),
two label-free alternatives (PAT~\cite{seth2024pat}, consistency-guided test-time adaptation~\cite{chen2024mctta}),
and the only paper I found that targets the same problem as my project, noise robustness of zero-shot audio-language classification (DAS~\cite{vo2026das}).
Inclusion criteria: (i) the paper introduces a language-audio model or adaptation method evaluated on ESC-50, or a mechanism I use; (ii) the full text is public.
I excluded music-only models and captioning/generation models to keep the review deep rather than wide
within the two-week budget. Eleven papers are analysed; numbers in the matrix are the authors' reported numbers, not mine.

\paragraph{Notation used in this report.}
\textbf{[Reported]} marks a number or claim from a paper or its official repository,
\textbf{[Measured]} marks a number I obtained by running the code in my repository, and
\textbf{[Interpretation]} marks my own reading of the evidence.

\begin{landscape}
\begin{table}
\centering
\scriptsize
\caption{Comparison matrix (one row per paper), part 1: the foundation model and earlier audio-text models. All accuracies are \textbf{[Reported]} by the authors. ZS = zero-shot, LP = linear probe.}
\label{tab:matrix}
\renewcommand{\arraystretch}{1.15}
\begin{tabularx}{\linewidth}{@{}>{\raggedright\arraybackslash}p{1.75cm}*{7}{>{\raggedright\arraybackslash}X}@{}}
\toprule
Paper & Question and contribution & Modalities and assumptions & Mechanism and objective & Data and protocol & Baselines and evidence & Resources and reproducibility & Limitations and connections \\
\midrule
CLIP~\cite{radford2021clip} (ICML 2021) &
Can natural-language supervision at web scale replace fixed label sets for vision? New: 400M-pair dataset and prompt-based zero-shot transfer. &
Image + text at training; at test only an image and class names turned into prompts. Assumes the test concepts are describable in text. &
Two encoders (ResNet/ViT, Transformer); symmetric InfoNCE over a batch with learnable temperature; zero-shot = nearest prompt embedding. &
WIT (400M pairs, not released); 30+ downstream datasets; prompt templates and ensembling; LP on frozen features. &
Strong: compares with supervised ImageNet models and LP baselines on many datasets, studies distribution shift. Zero-shot roughly matches a 4-shot LP on average. &
Weights and inference code released; training data and training code not. Largest models trained on 256--592 V100 GPUs for 12--18 days. &
Fails on fine-grained and abstract tasks; very few-shot LP can be worse than zero-shot. Template for every audio model below. \\
\midrule
Wav2CLIP~\cite{wu2022wav2clip} (ICASSP 2022) &
Can audio be put into CLIP space without any text? New: distil CLIP image embeddings into an audio encoder. &
Video (audio + frames) at training, no text; text enters only through the frozen CLIP text tower at test. &
ResNet-18 audio encoder trained with a contrastive loss to match the frozen CLIP image embedding of the same video. &
VGGSound ($\approx$200k 10\,s clips); evaluated on ESC-50, UrbanSound8K, etc.\ with ZS, LP and retrieval. &
ESC-50 ZS 41.4\%, LP 86.0\%. Good LP shows useful features; weak ZS shows image-mediated alignment is loose. &
Code and weights released; single, small encoder, cheap to run. &
Text alignment is indirect (via images), so ZS is the weakest in the table. First step towards CLIP-style audio. \\
\midrule
AudioCLIP~\cite{guzhov2022audioclip} (ICASSP 2022) &
Extend CLIP to a third modality. New: tri-modal (image, text, audio) training. &
Audio + image + class-label text at training; any pair at test. &
ESResNeXt audio head added to frozen/partially tuned CLIP; contrastive losses between all modality pairs. &
AudioSet (labels used as text); ESC-50 and UrbanSound8K, ZS and fine-tuned. &
ESC-50 ZS 69.4\%; fine-tuned 97.2\%. Ablations on which towers are trained. &
Code released. &
Text is only label names, so language is narrow. Bridge between Wav2CLIP and true audio-caption training. \\
\midrule
Microsoft CLAP~\cite{elizalde2023clap} (ICASSP 2023) &
Learn audio concepts directly from audio-caption pairs. New: CLIP recipe with audio and captions. &
Audio + free-text caption at training; audio + prompt at test. &
CNN14 audio encoder, BERT text encoder, linear projections, symmetric contrastive loss. &
128k pairs (FSD50K, Clotho, AudioCaps, MACS); 16 downstream tasks; prompt ``this is a sound of [label]''. &
ESC-50 ZS 82.6\%; prompt study shows template sensitivity. &
Code released; training uses public datasets. &
Training data shares the Freesound source with ESC-50; LAION-CLAP later lists FSD50K clips that overlap ESC-50, which makes leakage a real concern. \\
\bottomrule
\end{tabularx}
\end{table}
\begin{table}
\centering
\scriptsize
\caption*{\textbf{Table~\ref{tab:matrix} (continued).} The reproduced paper and the two few-shot adaptation papers.}
\begin{tabularx}{\linewidth}{@{}>{\raggedright\arraybackslash}p{1.75cm}*{7}{>{\raggedright\arraybackslash}X}@{}}
\toprule
Paper & Question and contribution & Modalities and assumptions & Mechanism and objective & Data and protocol & Baselines and evidence & Resources and reproducibility & Limitations and connections \\
\midrule
\textbf{LAION-CLAP}~\cite{wu2023laionclap} (ICASSP 2023), \textbf{reproduced} &
Scale CLAP and handle variable-length audio. New: LAION-Audio-630K, feature fusion, keyword-to-caption (K2C) augmentation. &
Audio (48\,kHz) + caption at training; for label-only data, captions are generated from labels by T5. Same zero-shot use as CLIP. &
HTSAT or PANN audio encoder, RoBERTa/BERT/CLIP text encoder, 2-layer MLP to 512-d, symmetric contrastive loss. Fusion: attention mix of a global down-sampled view and 3 local 10\,s views for audio $>$10\,s. &
633,526 pairs (4,325\,h) + AudioSet with K2C captions; overlapping clips removed from ESC-50 before ZS evaluation. Prompt ``This a sound of label.''. &
ESC-50 ZS: CLAP 89.1, +Fusion 88.0, +K2C 91.0\%. Mapping of these rows to released checkpoints is not stated. Fusion mainly helps long-audio retrieval (Clotho). &
Code, pip package and 4 general checkpoints released; dataset released as links only. GPU type/count not stated. &
Robustness to noisy input not evaluated; no few-shot adaptation. Strongest ZS in the table, so it is my baseline. \\
\midrule
Tip-Adapter~\cite{zhang2022tip} (ECCV 2022) &
Few-shot adaptation of CLIP without training. New: key-value cache of few-shot features added to zero-shot logits. &
Frozen CLIP; $K$ labelled images per class at adaptation time; no gradient updates (Tip-Adapter-F fine-tunes keys). &
$\text{logits} = f W^\top + \alpha\, \exp(-\beta(1 - f F^\top)) L$; $\alpha$ residual weight, $\beta$ sharpness. &
11 image datasets, 1--16 shots. &
Training-free version is competitive with CoOp and CLIP-Adapter, the fine-tuned version beats them at far lower training cost; ablations on $\alpha$, $\beta$, cache size. ImageNet default $\beta{=}5.5$, $\alpha{=}1$. &
Code released; runs in seconds on frozen features. &
Image only; cache built from clean shots; hyper-parameters need a validation set. Mechanism I transfer to audio. \\
\midrule
Treff adapter~\cite{liang2023treff} (arXiv 2023) &
Can a language-audio model be a few-shot learner? New: cross-attention linear model (CALM) on top of CLAP, with a training-free variant. &
Frozen CLAP; $K$ labelled clips per class; train-free or lightly trained adapter. &
$P = P_{\text{ZS}} + \alpha P_{\text{FS}}$, $P_{\text{FS}}$ from cross-attention between test and support embeddings. &
ESC-50, FSDKaggle2018, FSD-FS; 16-shot main setting. &
ESC-50 16-shot: CLAP ZS 82.6, training-free Treff 87.75, Tip-Adapter (tuned) 91.45, Treff tuned 92.21\%. &
Code ``to be released''; I could not find it. &
Built on a CLAP model with ESC-50 ZS 82.6 (the Microsoft CLAP figure); clean support and clean test only. Closest prior work to my improvement. \\
\bottomrule
\end{tabularx}
\end{table}
\begin{table}
\centering
\scriptsize
\caption*{\textbf{Table~\ref{tab:matrix} (continued).} Recent work that builds on CLAP-style models (2024--2026): competing adaptation methods and noise robustness.}
\begin{tabularx}{\linewidth}{@{}>{\raggedright\arraybackslash}p{1.75cm}*{7}{>{\raggedright\arraybackslash}X}@{}}
\toprule
Paper & Question and contribution & Modalities and assumptions & Mechanism and objective & Data and protocol & Baselines and evidence & Resources and reproducibility & Limitations and connections \\
\midrule
PALM~\cite{hanif2024palm} (EMNLP 2024) &
Does prompt learning (CoOp-style) work for audio-language models? New: learn a per-class offset in the text-feature space instead of input tokens. &
Frozen audio and text encoders; 16 labelled clips per class. &
$f'(t_i) = (1-\lambda_i) f_T(t_i) + \lambda_i z_i$ with learnable $z_i, \lambda_i$; cross-entropy on the few-shot set. &
Pengi backbone; 11 datasets incl.\ ESC-50 (5-fold); 3 seeds. &
ESC-50: zero-shot 49.65, CoOp 93.82, CoCoOp 94.27, PALM 95.93\%. No Tip-Adapter comparison. &
Project page and code released. &
Authors list noise robustness as future work. Competing labelled few-shot method to the cache I use; needs gradient training. \\
\midrule
PAT~\cite{seth2024pat} (arXiv 2024) &
Can zero-shot audio classification be improved without any training or labels? New: weighted prompt ensembles plus parameter-free audio-text attention. &
Frozen CLAP-style model; no labelled audio; two mixing weights chosen per dataset. &
$L = L_{\text{pred}} + \beta_1 L_{\text{audio}} + \beta_2 L_{\text{text}}$ from an attention map between frame-level audio and text features. &
6 backbones incl.\ LAION-CLAP and MS-CLAP; 18 datasets incl.\ ESC-50; robustness test with audio augmentations. &
ESC-50 LAION-CLAP 89.00 $\to$ 93.00\%; with Gaussian noise (MS-CLAP-23) 91.80 $\to$ 94.20\%. No SNR levels given. &
Code link given in a footnote. &
Noise is an augmentation test without controlled SNR; no labelled shots. Shows a label-free route to the same zero-shot gains. \\
\midrule
Consistency-guided TTA~\cite{chen2024mctta} (arXiv 2024) &
Can a CLAP model adapt at test time with unlabelled audio? New: context- and domain-aware prompt generators trained by consistency losses. &
Frozen CLAP; unlabelled test audio, small batches; no labels. &
Self-entropy over augmented views plus a contrastive consistency term across test samples. &
12 datasets incl.\ ESC-50; augmentations are spectrogram masking, not noise. &
ESC-50 zero-shot 93.90 $\to$ 94.65\%; 12-task average 62.93 $\to$ 68.83\%. &
Code not stated. &
Not evaluated under noise. Small gain on ESC-50 (+0.75), larger on other datasets; an unlabelled alternative to my labelled cache. \\
\midrule
DAS~\cite{vo2026das} (ICML 2026 audio workshop) &
Why does zero-shot CLAP fail under acoustic noise, and can text alone fix it? New: class-wise ``drift'' directions from noisy vs.\ clean prompts. &
Frozen LAION-CLAP; no labelled or unlabelled audio needed; text prompts only. &
$\text{score}(z,c) = z\cdot C_c + \beta\, z\cdot\hat\delta_c$, $\hat\delta_c$ = mean text embedding of ``$c$ with noise'' minus ``sound of $c$'', $\beta = 0.25$. &
Real urban scene noise (TAU 2019) at 0--20\,dB on UrbanSound8K and FSD50K; ESC-50 and other noises in the supplement. &
Accuracy falls 12--30 points at 0\,dB; UrbanSound8K at 0\,dB 54.80 $\to$ 59.15\%. ESC-50: mean +1.84 over 66 settings. &
Code not stated; training-free. &
Closest to my question, but zero-shot only: no labelled shots, no cache. My work is the labelled few-shot counterpart. \\
\bottomrule
\end{tabularx}
\end{table}
\end{landscape}

\paragraph{Synthesis.}
\textbf{[Interpretation]} The first seven papers form one line of development.
CLIP established the recipe (two towers, contrastive loss, prompts as classifiers).
Audio work first borrowed CLIP's space indirectly (Wav2CLIP through images, AudioCLIP through labels),
then trained on real audio captions (Microsoft CLAP), and LAION-CLAP scaled the data and added mechanisms for long audio and
label-only data. The ESC-50 zero-shot numbers climb from 41.4 to 91.0\%, but three issues limit how much of that climb can be attributed to method:
(1) training sets and ESC-50 share the Freesound source, and only LAION-CLAP documents overlap removal;
(2) prompt wording changes zero-shot accuracy, and the prompts differ between papers (LAION-CLAP prints a prompt with a typo);
(3) LAION-CLAP's table does not say which released checkpoint produced which row.
These motivate my reproduction design: several checkpoints, several prompts, all 2000 clips and the official fold-1 split.

On the adaptation side the recent work splits by what it needs at deployment.
\emph{Labelled few-shot} methods either keep the model frozen and add a cache (Tip-Adapter, Treff) or train a small set of parameters (PALM); both add several points with 16 clips per class.
\emph{Label-free} methods improve zero-shot scoring from text alone (PAT, DAS) or from unlabelled test audio (consistency-guided TTA).
Only DAS studies controlled background noise, and it reports the same order of degradation I measure (12--30 points at 0\,dB, Section~\ref{sec:results}); PAT's Gaussian-noise test has no stated SNR.

\textbf{Research gap.} Every labelled few-shot method in this review (Tip-Adapter, Treff, PALM) builds its support set from clean recordings and is tested on clean recordings,
and the noise-robustness work (DAS, PAT) uses no labelled clips at all.
No paper asks whether the few labelled clips themselves can be used to absorb noise.
A cache is a natural place to do this: keys are just embeddings, so adding embeddings of noisy versions of the same support clips costs one extra forward pass per clip and no training.
This is the question of Section~\ref{sec:improve}. The label-free methods are complementary rather than competing: DAS's text-derived drift could be combined with a noisy cache, which I leave open.
